\begin{resumo}[Abstract]
\begin{otherlanguage*}{english}

The practical teaching of Requirements Engineering (RE) faces significant logistical challenges in simulating the ambiguities and sociotechnical complexity inherent in the real elicitation process. Although the ElicitBot conversational agent proof of concept demonstrated initial pedagogical potential in individual sessions, limitations regarding architectural coupling, in-memory state volatility, and cumulative growth in token consumption prevented its simultaneous operation in full classes. This work aimed to technically and functionally enable the ElicitBot tool, allowing its controlled use as a pedagogical resource to support requirements elicitation practice in Software Requirements classes. To this end, a methodology guided by the principles of Design Science Research (DSR) was adopted, structured into four phases: diagnosis and planning, refactoring and development, field deployment and validation, and evaluation and results analysis. Based on the diagnosis of the inherited system, the elaboration of eight Architecture Decision Records (ADRs), and financial modeling grounded in Cloud Economics, the application was restructured into a decoupled client-server architecture (React, FastAPI, and PostgreSQL) under the Read-Release-Write transactional pattern and deployed on a containerized institutional server. The conversational engine evolved into a 12-stage Finite State Machine coupled with bipartite memory management using Context Caching, lifecycle control of eight traceable artifact categories supporting Human-in-the-Loop intervention, and dual validation of Functional Requirements, achieving a 0\% failure rate for up to 50 concurrent users in load and soak tests conducted with the Locust tool. Empirical classroom validation involved 78 students distributed across 14 teams simultaneously (processing 519 conversational turns), the collection of $N = 65$ valid individual evaluations ($\alpha = 0.877$), and the immediate application of corrective refinements to the interface and system prompt. Results showed a median and mode equal to the maximum score (5.0) across all 20 evaluated criteria (overall mean of $\mu = 4.49$), highlighted by 100.0\% approval in logical flow adherence, 96.6\% in requirements traceability, 90.8\% satisfaction with the interaction, and 83.1\% adoption potential, alongside a 74.0\% cache hit rate across turns, a median latency of 5.14 s, and a total billed cost of R\$ 20.84 to serve both evaluated classes. It is concluded that combining persistent asynchronous infrastructure, state-driven context engineering, and institutional governance and user control mechanisms enables the scalable and economically sustainable adoption of Large Language Model (LLM) agents to support practical Requirements Engineering education.

\vspace{\onelineskip}
\noindent
\textbf{Keywords}: Requirements Engineering. Artificial Intelligence. Large Language Models. Smart Education. ElicitBot.
\end{otherlanguage*}
\end{resumo}
